\section*{अध्याय \ref{ch:b3:chromatin-epigenetics} --- क्रोमैटिन और एपिजेनेटिकी}

\begin{solution}{exo:b3:chromatin-epigenetics:1}
H2A, H2B, H3 और H4 की दो-दो प्रतियाँ (एक अष्टलक), जिनके साथ
$1.65$ फेरों में \qty{147}{bp} DNA; \qtyrange{20}{60}{bp} के योजक,
जिनसे लगभग \qty{200}{bp} की पुनरावृत्ति बनती है। न्यूक्लिएज़ केवल खुला
योजक DNA काटती है, किसी क्रोड के भीतर कभी नहीं, इसलिए हर खंड
पुनरावृत्तियों की पूर्ण संख्या है: \num{200}, \num{400},
\qty{600}{bp} की सीढ़ी, धब्बा नहीं।
\end{solution}

\begin{solution}{exo:b3:chromatin-epigenetics:2}
द्विगुणित: $6.0\times 10^{9}$ bp; $6.0\times 10^{9}/190 \approx
3.2\times 10^{7}$ \omterm{def:b3:chromatin-epigenetics:nucleosome}{न्यूक्लियोसोम}; लंबाई $6.0\times 10^{9}\times
\qty{0.34}{nm} = \qty{2.0}{m}$।
\end{solution}

\begin{solution}{exo:b3:chromatin-epigenetics:3}
H3K27me3: \omterm{def:b3:chromatin-epigenetics:nucleosome}{हिस्टोन} H3, लाइसीन 27, त्रिमेथिलित --- परिवर्धन में दमित
जीनों का मौन \omterm{def:b3:chromatin-epigenetics:nucleosome}{क्रोमैटिन}, जिसे PRC2 लिखती है और PRC1 के क्रोमोप्रांत
प्रोटीन पढ़ते हैं। H4K16ac: H4, लाइसीन 16, ऐसीटिलित --- सक्रिय, खुला
\omterm{def:b3:chromatin-epigenetics:nucleosome}{क्रोमैटिन} (यह एक न्यूक्लियोसोम--न्यूक्लियोसोम संपर्क तोड़ता है)।
H3S10ph: H3, सेरीन 10, फ़ॉस्फ़ोरिलित --- गुणसूत्र संघनन से जुड़ा एक
समसूत्री चिह्न (और, क्षणिक रूप से, कुछ संकेतों से जीन-सक्रियण से भी);
कोई स्थायी मौनकारी या सक्रियकारी चिह्न नहीं।
\end{solution}

\begin{solution}{exo:b3:chromatin-epigenetics:4}
नारंगी जीन X-सहलग्न है और उसके एक नारंगी तथा एक काला युग्मविकल्पी हैं;
धब्बेदार रोम के लिए दो भिन्न X गुणसूत्रों पर हर एक का होना चाहिए,
जिनमें से प्रति कोशिका एक मौन कर दिया जाता है। किसी नर के पास एक ही X
है और वह पूरा नारंगी या पूरा काला होता है। कोई कैलिको नर दो X गुणसूत्र
और एक Y (XXY) ढोता है, और बंध्य होता है।
\end{solution}

\begin{solution}{exo:b3:chromatin-epigenetics:5}
प्रति \qty{11}{nm} मनके पर
$\qty{200}{bp}\times \qty{0.34}{nm} = \qty{68}{nm}$: अनुपात $6.2$।
आवश्यक: $\qty{4}{cm}/\qty{4}{\micro m} =
10^{4}$। और मोड़ना $10^{4}/6.2 \approx 1600$ गुना।
\end{solution}

\begin{solution}{exo:b3:chromatin-epigenetics:6}
$\mu = 0.98$, $\delta = 0.01$: $p^{*} = 0.01/0.03 = 0.33$; अनुपात $\mu -
\delta = 0.97$; अर्ध-आयु $\ln 2/(-\ln 0.97) = 22.8 \approx 23$
विभाजन। $\mu = 0.999$, $\delta = 0.001$: $p^{*} = 0.5$; अनुपात
$0.998$; अर्ध-आयु $\ln 2 / 0.002 \approx 350$ विभाजन।
\end{solution}

\begin{solution}{exo:b3:chromatin-epigenetics:7}
$A$: सक्रिय (प्रवर्तक पर H3K4me3 के साथ \omterm{def:b3:chromatin-epigenetics:histone-code}{ऐसीटिलीकरण})। $B$: मौन,
पॉलीकोम्ब-दमित --- ऐच्छिक \omterm{def:b3:chromatin-epigenetics:nucleosome}{विषमक्रोमैटिन}, यानी ऐसे परिवर्धन-जीन की
अवस्था जो इस वंशक्रम में बंद है पर दूसरों में काम आता है, इसलिए किसी
दूसरे कोशिका-प्रकार में सक्रिय होने की संभावना $B$ की ही है। $C$: मौन,
संघटक \omterm{def:b3:chromatin-epigenetics:nucleosome}{विषमक्रोमैटिन} (HP1 पथ), प्रायः पुनरावृत्तियाँ या हर
कोशिका-प्रकार में मौन जीन।
\end{solution}

\begin{solution}{exo:b3:chromatin-epigenetics:8}
\emph{UBE3A} तंत्रिका-कोशिकाओं में केवल मातृ युग्मविकल्पी से अभिव्यक्त
होता है। उसका विलोपन उसके पिता से आया, उस युग्मविकल्पी पर जो वैसे भी
मौन है, इसलिए वह स्वस्थ है; जब वह उसे आगे देती है तो वह मातृ विलोपन बन
जाता है: हर बच्चे को उसे पाने की प्रायिकता $1/2$ है और तब उसे एंजेलमैन
संलक्षण होता है। यदि विलोपन उसकी माता से आया होता तो उसे स्वयं एंजेलमैन
संलक्षण होता; कोई स्वस्थ वाहक उसे केवल पैतृक रूप से ही पा सकता है।
(पूरे क्षेत्र का पैतृक विलोपन प्रैडर--विली देगा; अकेला \emph{UBE3A}
उसके लिए पर्याप्त नहीं।)
\end{solution}

\begin{solution}{exo:b3:chromatin-epigenetics:9}
जिस X पर \emph{\omterm{def:b3:chromatin-epigenetics:x-inactivation}{Xist}} नहीं है वह निष्क्रिय नहीं हो सकता, इसलिए हर कोशिका
में दूसरा X मौन होता है: निष्क्रियण पूरा होता है पर अब यादृच्छिक नहीं
(तिरछा), और भ्रूण जीवक्षम है। दोनों X पर \emph{\omterm{def:b3:chromatin-epigenetics:x-inactivation}{Xist}} न हो: कोई
निष्क्रियण नहीं, दो सक्रिय X गुणसूत्र, दोगुनी X-सहलग्न अभिव्यक्ति, और
मादा भ्रूण जल्दी मर जाता है। किसी अलिंगसूत्र पर \emph{\omterm{def:b3:chromatin-epigenetics:x-inactivation}{Xist}} पारजीन उस
अलिंगसूत्र को सिस में ढक लेता है और उसके जीन मौन कर देता है, जिससे एक
अलिंगसूत्री प्रति का प्रकार्य-ह्रास होता है --- यही वह प्रयोग है जिसने
दिखाया कि \emph{\omterm{def:b3:chromatin-epigenetics:x-inactivation}{Xist}} सिस में मौनीकरण के लिए पर्याप्त है।
\end{solution}

\begin{solution}{exo:b3:chromatin-epigenetics:10}
मेथिलित साइटोसीन विअमीनित होकर थाइमीन बनते हैं, जिसे मरम्मत पहचानती
नहीं, इसलिए कोई मेथिलित \omterm{def:b3:chromatin-epigenetics:methylation}{CpG} उच्च दर से TpG या CpA में उत्परिवर्तित
होता है। केवल वे अनुक्रम जो पीढ़ी-दर-पीढ़ी \emph{जनन-वंश} में अमेथिलित
रहे हैं अपने \omterm{def:b3:chromatin-epigenetics:methylation}{CpG} बचाए रखते हैं; द्वीपों ने अपने बचाए हैं, इसलिए वे
कशेरुकी इतिहास के अधिकांश भाग में जनन-वंश में अमेथिलित रहे हैं। केवल
यकृत में मेथिलित कोई जीन शुक्राणु और अंडाणु में अमेथिलित है; यकृत
कोशिकाओं में उठे उत्परिवर्तन वंशागत नहीं होते, इसलिए कायिक मेथिलीकरण
अनुक्रम में कोई पदचिह्न नहीं छोड़ता।
\end{solution}

\begin{solution}{exo:b3:chromatin-epigenetics:11}
प्रतिकृति पर पुराने, चिह्नित \omterm{def:b3:chromatin-epigenetics:nucleosome}{हिस्टोन} दोनों पुत्री द्विकुंडलियों में
बँट जाते हैं और नए, अचिह्नित \omterm{def:b3:chromatin-epigenetics:nucleosome}{हिस्टोनों} के बीच-बीच में बैठते हैं, इसलिए
हर पुत्री प्रांत लगभग आधा चिह्नित होता है। बचे हुए किसी चिह्नित
\omterm{def:b3:chromatin-epigenetics:nucleosome}{न्यूक्लियोसोम} से बँधा HP1 SUV39H को बुलाता है, जो अचिह्नित पड़ोसियों का
मेथिलीकरण करती है; जब तक चिह्नित \omterm{def:b3:chromatin-epigenetics:nucleosome}{न्यूक्लियोसोम} इतने सघन हैं कि हर
अंतराल तक कोई \omterm{def:b3:chromatin-epigenetics:histone-code}{लेखक} ला सकें, प्रांत अगले विभाजन से पहले फिर भर जाता है
--- वही पुनर्भरण जो प्रमेय का ऊँचा $\mu$ है। फैलाव सीमाओं पर रुकता है:
विद्युतरोधक प्रोटीन (CTCF), अत्यधिक अनुलेखित जीन और tRNA जीन, ऐसे
\omterm{def:b3:chromatin-epigenetics:nucleosome}{न्यूक्लियोसोम} जिन पर असंगत सक्रिय चिह्न (H3K4me3, \omterm{def:b3:chromatin-epigenetics:histone-code}{ऐसीटिलीकरण}) हैं
जिन्हें \omterm{def:b3:chromatin-epigenetics:histone-code}{लेखक} उपयोग नहीं कर सकती, और वह बिंदु जहाँ \omterm{def:b3:chromatin-epigenetics:histone-code}{विलोपक} (डीमेथिलेज़,
HDAC का आवागमन) \omterm{def:b3:chromatin-epigenetics:histone-code}{लेखक} से आगे निकल जाते हैं। परीक्षण: HP1 को किसी
प्रतिवेदक पर क्षणिक रूप से बुलाइए (ऐसे DNA-बंधक प्रांत के साथ संलयन
जिसका बंधन औषध से नियंत्रित हो), फिर उसे छोड़ दीजिए, और विभाजनों के
पार प्रतिवेदक की मौनता तथा H3K9me3 का पीछा कीजिए; पूर्वानुमान यह है कि
बने रहना \omterm{def:b3:chromatin-epigenetics:histone-code}{लेखक} पर निर्भर होगा, और \omterm{def:b3:chromatin-epigenetics:histone-code}{लेखक} या \omterm{def:b3:chromatin-epigenetics:histone-code}{पाठक} के द्विलकन के
उत्परिवर्तित होने पर, या प्रांत के भीतर कोई विद्युतरोधक रख देने पर,
वह खो जाएगा।
\end{solution}

\begin{solution}{exo:b3:chromatin-epigenetics:12}
हटाइए: साझा परिवेश (पोते-पोतियों का अपना आहार और निर्धनता); गर्भाशयी
प्रभाव --- अकाल-संपर्क में आई माता की पुत्री तब भ्रूण थी और उसकी
जनन-कोशिकाएँ पहले से मौजूद थीं, इसलिए पोते का जनन-वंश सीधे संपर्क में
आया और प्रभाव पीढ़ी-पार नहीं है; उन परिवारों के बीच आनुवंशिक भेद जो बचे
और जो नहीं बचे; उलटा कारणत्व (शरीर-द्रव्यमान स्वयं रक्त में
मेथिलीकरण बदलता है); रक्त-कोशिका संघटन के भेद; और अनेक \omterm{def:b3:chromatin-epigenetics:methylation}{CpG} पर बहु
परीक्षण। चूहों में प्रयोग: अंतःप्रजात F$_0$ \emph{नरों} को संपर्क में
लाइए (कोई गर्भाशयी मार्ग नहीं), उनका संगम असंपर्कित मादाओं से कराइए या
पात्रे निषेचन तथा भ्रूण-स्थानांतरण का उपयोग कीजिए, और पैतृक वंश से
F$_2$ तथा F$_3$ का पीछा कीजिए, हर पीढ़ी के शुक्राणु में मेथिलीकरण
मापते हुए; बिना किसी आनुवंशिक विचरण वाली प्रजाति में जो चिह्न F$_3$ के
शुक्राणु और लक्षण-प्ररूप में बना रहे, वह जनन-वंशीय संचरण है।
\end{solution}

\begin{solution}{pb:b3:chromatin-epigenetics:1}
\textbf{1.} $6.4\times 10^{9}\times \qty{0.34}{nm} = \qty{2.2}{m}$;
$6.4\times 10^{9}/200 = 3.2\times 10^{7}$ \omterm{def:b3:chromatin-epigenetics:nucleosome}{न्यूक्लियोसोम}।
\textbf{2.} \omterm{def:b3:chromatin-epigenetics:nucleosome}{हिस्टोन}: $3.2\times 10^{7}\times 1.08\times 10^{5} =
3.5\times 10^{12}$ Da $= \qty{5.7}{pg}$। DNA: $6.4\times 10^{9}\times
650 = 4.2\times 10^{12}$ Da $= \qty{6.9}{pg}$। तुलनीय द्रव्यमान।
\textbf{3.} केंद्रक: $\tfrac{4}{3}\pi\,(\qty{5}{\micro m})^{3} =
\qty{520}{\micro m^3}$। क्रोड: $\pi\,(5.5)^{2}\times 5.5 =
\qty{520}{nm^3} = 5.2\times 10^{-7}\,\unit{\micro m^3}$; कुल
$3.2\times 10^{7}\times 5.2\times 10^{-7} = \qty{17}{\micro m^3}$:
केंद्रक का लगभग \qty{3}{\%}।
\textbf{4.} $3.2\times 10^{7}\times \qty{11}{nm} = \qty{0.35}{m}$;
\omterm{prop:b3:chromatin-epigenetics:compaction}{संवेष्टन अनुपात} $2.2/0.35 \approx 6$।
\textbf{5.} $6.4\times 10^{9}/2\times 10^{5} = \num{32000}$ पाश,
प्रत्येक \num{1000} \omterm{def:b3:chromatin-epigenetics:nucleosome}{न्यूक्लियोसोमों} का।
\textbf{6.} $0.21^{2}\times 6.4\times 10^{9} = 2.8\times 10^{8}$;
प्रेक्षित $5.6\times 10^{7}$, अर्थात् \qty{80}{\%} खो गए,
5-मेथिलसाइटोसीन के थाइमीन में विअमीनन से (मेथिलित \omterm{def:b3:chromatin-epigenetics:methylation}{CpG} पर बिना मरम्मत
हुए C$\to$T संक्रमण)।
\textbf{7.} $p_{n+1} = 0.99\,p_{n} + 0.02\,(1 - p_{n}) = 0.97\,p_{n} +
0.02$; $p^{*} = 0.02/0.03 = 0.67$।
\textbf{8.} $p_{n} = 0.67 + 0.33\times 0.97^{n}$: $n = 10$: $0.91$;
$n = 50$: $0.74$; $n = 200$: $0.67$ (आधिक्य $7\times 10^{-4}$ है)।
\textbf{9.} $\ln 2/(-\ln 0.97) = 22.8 \approx 23$ विभाजन; सप्ताह में
एक की दर से लगभग \num{23} सप्ताह --- पाँच महीने।
\textbf{10.} हर प्रतिकृति एक पुरानी, अब भी मेथिलित लड़ी को एक ऐसी नई
लड़ी से जोड़ती है जिसका मेथिलीकरण कोई एंज़ाइम नहीं करती; पुरानी लड़ियाँ
न बढ़ती हैं न घटतीं, लड़ियों की संख्या दोगुनी होती है, इसलिए मेथिलित
अंश आधा हो जाता है। $0.80/2^{n} < 0.05$ के लिए $2^{n} > 16$ चाहिए:
पाँच विभाजन ($0.025$; चार से ठीक \qty{5}{\%} मिलता है)।
\textbf{11.} अनुपात $0.9995 - 0.0005 = 0.999$; अर्ध-आयु $\ln
2/(-\ln 0.999) \approx 690$ विभाजन, सप्ताह में एक की दर से लगभग
\num{13} वर्ष।
\textbf{12.} जीवन भर के लगभग \num{4000} साप्ताहिक विभाजनों तक मौन
रहने वाले किसी जीन को सैकड़ों विभाजनों की अर्ध-आयु चाहिए, जो केवल
प्रश्न~11 का पुनर्भरण देता है: चिह्न के \omterm{def:b3:chromatin-epigenetics:histone-code}{पाठक} (MeCP2, HP1) उसके \omterm{def:b3:chromatin-epigenetics:histone-code}{लेखक}
(H3K9 मेथिलेज़, DNMT3) बुलाते हैं, इसलिए चिह्नों का कोई सघन प्रांत ऐसी
दक्षता से अपनी नकल करता है जो किसी अकेले एंज़ाइम के पास नहीं।
\textbf{13.} एक X पर नारंगी युग्मविकल्पी है, दूसरे पर काला; हर
त्वचा-कोशिका ने एक X मौन कर दिया है और केवल दूसरा अभिव्यक्त करती है;
किसी कोशिका के वंशज उसका चुनाव बनाए रखते हैं, इसलिए कोई धब्बा एक क्लोन
है (या एक ही चुनाव करने वाले सटे हुए क्लोनों का समूह)।
\textbf{14.} सभी $N$ नारंगी वाला X चुनें, या सभी दूसरा चुनें:
$2\times (1/2)^{N} = 2^{1-N}$।
\textbf{15.} $2^{1-N} = 10^{-9}$: $N - 1 = \log_{2} 10^{9} = 29.9$,
$N \approx 30$ कोशिकाएँ।
\textbf{16.} $\qty{3000}{cm^2}/\qty{15}{cm^2} = 200$ धब्बे, जो $30$
संस्थापकों से कहीं अधिक हैं: वृद्धि के दौरान किसी संस्थापक के वंशज
कोशिका-मिश्रण और प्रवसन से बिखर जाते हैं, इसलिए एक क्लोन कई द्वीप
बनाता है (और एक ही चुनाव वाले पड़ोसी आपस में मिल जाते हैं, जो उलटी
दिशा में काम करता है पर कम)।
\textbf{17.} पहले X के बाद हर X पर एक \omterm{def:b3:chromatin-epigenetics:x-inactivation}{बार पिंड}: XXY में एक, XXX में
दो, XY में कोई नहीं।
\textbf{18.} सहनीय तब, जब जीन का कोई प्रकार्यशील Y समजात हो, ताकि नरों
के पास भी दो प्रतियाँ हों (छद्म-अलिंगसूत्री जीन और
\emph{KDM5C}/\emph{KDM5D} जैसे युग्म); यह X असुगुणिताओं में मायने रखता
है --- टर्नर संलक्षण (XO) में बच निकलने वाले जीन एक ही प्रति में होते
हैं, इसीलिए कोई XO व्यक्ति बस कोई सामान्य मादा नहीं है (\emph{SHOX}
की अर्धअपर्याप्तता और छोटा कद), और XXY में बच निकलने वालों की मात्रा
अधिक हो जाती है।
\textbf{19.} $1/\num{15000} = 6.7\times 10^{-5}$; विलोपन $0.70\times
6.7\times 10^{-5} = 4.7\times 10^{-5}$ (\num{21000} में एक); मातृ
द्विसूत्रता $1.7\times 10^{-5}$ (\num{60000} में एक)।
\textbf{20.} एंजेलमैन एक ही जीन, \emph{UBE3A}, का ह्रास है, इसलिए मातृ
युग्मविकल्पी में एक ही बिंदु उत्परिवर्तन पर्याप्त है। प्रैडर--विली एक
साथ कई पैतृक रूप से अभिव्यक्त जीनों का ह्रास है (\emph{SNRPN},
लघु-RNA गुच्छा); कोई बिंदु उत्परिवर्तन उनमें से एक को निष्क्रिय करता
है और अधिक से अधिक आंशिक लक्षण-प्ररूप देता है, कभी पूरा संलक्षण नहीं।
\textbf{21.} \qty{20}{\%} बच्चों को पैतृक विलोपन मिलता है और उन्हें
प्रैडर--विली होता है; किसी को एंजेलमैन नहीं, जिसके लिए मातृ विलोपन
चाहिए।
\textbf{22.} तीन गुणसूत्र, दो मातृ और एक पैतृक, एक यादृच्छिक रूप से
खोया जाता है: पैतृक प्रायिकता $1/3$ से खोता है, जिससे मातृ
द्विसूत्रता और प्रैडर--विली बनते हैं; प्रायिकता $2/3$ से कोई मातृ
प्रति खोती है और बच्चा सामान्य होता है।
\textbf{23.} पैतृक जीनों को माता पर माँग बढ़ानी चाहिए; उनका ह्रास
(प्रैडर--विली) माँग घटाएगा --- जो नवजात के दुर्बल स्तनपान और न पनपने
से मेल खाता है --- जबकि बाद की अतृप्त भूख इसका उलटा है और उसे मिलाना
कठिन है (एक प्रस्तावित पाठ यह है कि पैतृक जीन सामान्यतः बच्चे को
स्तन-त्याग की ओर मोड़ते हैं)। एंजेलमैन, किसी मातृ जीन का ह्रास, माँग
बढ़ाएगा; एंजेलमैन बच्चों का लंबा चूसना और बार-बार रात में जागना मेल
खाते हैं, दौरे और वाणी-ह्रास संसाधनों के बारे में हैं ही नहीं।
\textbf{24.} पैतृक प्रतिसंवेदी अनुलेख (\emph{UBE3A-ATS}) को ऐसे
प्रतिसंवेदी ऑलिगोन्यूक्लियोटाइड से लक्ष्य कीजिए जो उसे अपघटित कर दे
या उसका दीर्घन रोक दे: पैतृक \emph{UBE3A} अक्षत है और केवल सिस में उसी
RNA से मौन किया गया है, इसलिए RNA हटाने से उसकी अभिव्यक्ति लौट आती है।
\omterm{def:b3:chromatin-epigenetics:imprinting}{अंकन नियंत्रण क्षेत्र} का मेथिलीकरण अछूता रहता है, इसलिए \emph{SNRPN}
और दूसरे अंकित जीन अपना सामान्य जनकीय प्रतिरूप बनाए रखते हैं।
\textbf{25.} किसी चिह्न की अर्ध-आयु: बिना पुनर्भरण लगभग $23$ विभाजन,
उसके साथ लगभग $690$; कैलिको रोम का निर्णय लगभग $30$ कोशिकाओं की
संस्थापक समष्टि ने किया।
\end{solution}
